\documentclass[11pt]{article}
\usepackage[margin=1in]{geometry}
\usepackage[T1]{fontenc}
\usepackage{lmodern}
\usepackage{amsmath,amssymb,amsthm}
\usepackage{graphicx}
\usepackage[hidelinks]{hyperref}
\usepackage{microtype}
\emergencystretch=2em

\newtheorem{theorem}{Theorem}
\newtheorem{remark}{Remark}

\newcommand{\K}{\mathbb K}
\newcommand{\N}{\mathbb N}
\newcommand{\Pmeas}{\mathcal P}

\title{Signed uniformly distributed sequences from ordinary ones}
\author{Run fa\_banach\_001, agent\_super\_00}
\date{2026-06-30}

\begin{document}
\maketitle

\section*{Source Question}

The source paper is Mercourakis and Vassiliadis,
\emph{Uniform Distribution of Sequences and its interplay with Functional
Analysis}, arXiv:2102.02306.  Definition 1 of that paper says that, for a
compact Hausdorff space \(K\) and a real signed regular Borel measure
\(\mu\in M(K)\) with \(\|\mu\|=|\mu|(K)=1\), a sequence
\((x_n)\subseteq K\) is \(\mu\)-u.d. if
\begin{enumerate}
\item \((x_n)\) is \(|\mu|\)-uniformly distributed in the classical sense, and
\item there are signs \(\varepsilon_n\in\{\pm1\}\) such that
\[
  \frac{1}{N}\sum_{n=1}^N \varepsilon_n\delta_{x_n}
       \overset{w^*}{\longrightarrow}\mu .
\]
\end{enumerate}
The source proves this property for compact metric \(K\).  In concluding
future-work item 2(b), it asks about compact non-metrizable spaces \(K\) for
which every probability measure \(\nu\in P(K)\) admits an ordinary u.d.
sequence:
\begin{quote}
For such spaces it would be interesting to know if every signed measure
\(\mu\) with \(\|\mu\|=1\) admits a u.d. sequence in the sense of Definition 1.
\end{quote}

\begin{center}
\includegraphics[width=0.95\linewidth]{figures/open_problem_crop.png}
\end{center}

\section*{Answer}

\begin{theorem}\label{thm:main}
Let \(K\) be a compact Hausdorff space with the following property:
every probability measure \(\nu\in P(K)\) admits a uniformly distributed
sequence in the classical sense.  Then every real signed regular Borel measure
\(\mu\in M(K)\) with \(\|\mu\|=1\) admits a \(\mu\)-u.d. sequence in the sense
of Definition 1 of the source paper.
\end{theorem}

\section*{Proof Intuition}

The only obstruction in the signed setting is that the same points must
simultaneously realize the total variation \(|\mu|\) and, after signs are
attached, realize \(\mu\).  The Jordan decomposition separates these two
tasks.  Choose an ordinary uniformly distributed sequence for the normalized
positive variation and another for the normalized negative variation.  Then
merge the two lists so that their asymptotic proportions are exactly the two
Jordan masses.  Testing against any \(f\in C(K)\), the unsigned Cesaro
averages converge to the weighted sum \(\mu^++\mu^-=|\mu|\), while the signed
averages converge to \(\mu^+-\mu^-=\mu\).

\section*{Proof}

Let \(\mu=\mu^+-\mu^-\) be the Jordan decomposition of \(\mu\), and set
\[
   \alpha=\mu^+(K),\qquad \beta=\mu^-(K).
\]
Since \(\|\mu\|=|\mu|(K)=\mu^+(K)+\mu^-(K)\), we have
\[
   \alpha+\beta=1.
\]

If \(\alpha=1\), then \(\mu=|\mu|\) is a probability measure.  By the
hypothesis on \(K\), choose a \(\mu\)-u.d. sequence in the ordinary sense and
take all signs equal to \(+1\).  This is a \(\mu\)-u.d. sequence in the
source's signed sense.  If \(\beta=1\), choose an ordinary \(|\mu|\)-u.d.
sequence and take all signs equal to \(-1\).  Thus assume from now on that
\(\alpha,\beta>0\).

By the hypothesis on \(K\), choose sequences \((a_i)\subseteq K\) and
\((b_j)\subseteq K\) such that
\[
 \frac{1}{m}\sum_{i=1}^m \delta_{a_i}\overset{w^*}{\longrightarrow}
       \frac{\mu^+}{\alpha},
 \qquad
 \frac{1}{m}\sum_{j=1}^m \delta_{b_j}\overset{w^*}{\longrightarrow}
       \frac{\mu^-}{\beta}.
\]
Choose a subset \(A\subseteq\N\) with asymptotic density \(\alpha\).  For
example,
\[
   A=\{n\in\N:\lfloor n\alpha\rfloor>\lfloor(n-1)\alpha\rfloor\}.
\]
Then, with \(A_N=A\cap\{1,\ldots,N\}\), one has
\[
   |A_N|=\lfloor N\alpha\rfloor,\qquad
   \frac{|A_N|}{N}\to\alpha.
\]
Let \(B=\N\setminus A\); then \(|B\cap\{1,\ldots,N\}|/N\to\beta\).

Define a sequence \((x_n)\subseteq K\) and signs \((\varepsilon_n)\) as
follows.  If \(n\) is the \(i\)-th element of \(A\), put
\[
   x_n=a_i,\qquad \varepsilon_n=+1.
\]
If \(n\) is the \(j\)-th element of \(B\), put
\[
   x_n=b_j,\qquad \varepsilon_n=-1.
\]

Fix \(f\in C(K)\).  Write \(p_N=|A_N|\) and \(q_N=N-p_N\).  Since
\(\alpha,\beta>0\), both \(p_N\) and \(q_N\) tend to infinity.  Therefore
\[
\begin{aligned}
 \frac{1}{N}\sum_{n=1}^N f(x_n)
 &= \frac{p_N}{N}\left(\frac{1}{p_N}\sum_{i=1}^{p_N}f(a_i)\right)
    +\frac{q_N}{N}\left(\frac{1}{q_N}\sum_{j=1}^{q_N}f(b_j)\right)\\
 &\longrightarrow
   \alpha\int_K f\,d\frac{\mu^+}{\alpha}
  +\beta\int_K f\,d\frac{\mu^-}{\beta}
  =\int_K f\,d|\mu|.
\end{aligned}
\]
Thus \((x_n)\) is \(|\mu|\)-u.d. in the classical sense.

Similarly,
\[
\begin{aligned}
 \frac{1}{N}\sum_{n=1}^N \varepsilon_n f(x_n)
 &= \frac{p_N}{N}\left(\frac{1}{p_N}\sum_{i=1}^{p_N}f(a_i)\right)
    -\frac{q_N}{N}\left(\frac{1}{q_N}\sum_{j=1}^{q_N}f(b_j)\right)\\
 &\longrightarrow
   \alpha\int_K f\,d\frac{\mu^+}{\alpha}
  -\beta\int_K f\,d\frac{\mu^-}{\beta}
  =\int_K f\,d\mu.
\end{aligned}
\]
Equivalently,
\[
   \frac{1}{N}\sum_{n=1}^N \varepsilon_n\delta_{x_n}
      \overset{w^*}{\longrightarrow}\mu .
\]
The two defining conditions from Definition 1 are satisfied, so
\((x_n)\) is \(\mu\)-u.d.  This proves the theorem. \(\square\)

\section*{Consequences for the Source Question}

Theorem~\ref{thm:main} answers future-work item 2(b) affirmatively for every
compact Hausdorff space \(K\) satisfying the probability-measure hypothesis
stated in that item.  In particular, whenever an external theorem gives that
all probability measures on one of the non-metrizable compact spaces mentioned
there admit ordinary u.d. sequences, the signed-measure conclusion follows
immediately by the theorem above.

The argument does not address the separate future-work item 2(a), concerning
functions of bounded variation of several variables, nor the locally compact
noncompact direction in item 2(c).

\section*{Verification Notes}

No numerical or computer-assisted verification is involved.  The proof reduces
the signed problem to two ordinary probability-measure u.d. sequences and a
deterministic density partition of \(\N\).  The weak-star convergence is
checked directly on arbitrary \(f\in C(K)\).

\section*{Novelty Check}

Before promotion, the run's cheap indexes were searched for arXiv:2102.02306,
the title, the authors, and core phrases including ``signed measure u.d.
sequence'', ``signed measure uniformly distributed'', and ``compact
non-metrizable signed measure''.  No prior packet, attempt, or proof-gap entry
for this target was found.

Local corpus searches found the source paper and unrelated signed-sum/vector
sequence material, but no later solution of the future-work item.  Web/arXiv
spot checks on 2026-06-30 for exact phrases around ``2102.02306 signed measure
uniformly distributed'', ``every signed measure admits a u.d. sequence'',
``compact non-metrizable signed measure u.d. sequence'', and ``Mercourakis
Vassiliadis Uniform Distribution of Sequences'' found the source paper but no
exact later resolution.

\section*{Human Review Recommendation}

The main review point is whether the source future-work item 2(b) is intended
exactly under the hypothesis that every probability measure on \(K\) admits an
ordinary u.d. sequence.  Under that reading, the proof is a complete
affirmative answer.  The density interleaving calculation is the only
substantive step to check.

\begin{thebibliography}{9}

\bibitem{MercourakisVassiliadis}
S. K. Mercourakis and G. Vassiliadis,
\emph{Uniform Distribution of Sequences and its interplay with Functional
Analysis}, arXiv:2102.02306.

\end{thebibliography}

\end{document}
